% ============================================================
% DOCUMENTCLASS
% ============================================================
\documentclass[12pt,oneside]{book}

% ============================================================
% METADATOS
% ============================================================
\newcommand{\universidad}{Universidad de Buenos Aires}
\newcommand{\facultad}{Facultad de Derecho}
\newcommand{\materia}{Constitucional Profundizado}
\newcommand{\titulo}{Defensa del Consumidor y Defensa de la Competencia}
\newcommand{\subtitulo}{Concentraciones económicas, conductas anticompetitivas y garantías del control administrativo}
\newcommand{\autor}{Isaac Edgar Camacho Ocampo}
\newcommand{\anio}{2026}

% ============================================================
% PAQUETES
% ============================================================
\usepackage[utf8]{inputenc}
\usepackage[T1]{fontenc}
\usepackage[spanish,provide=*]{babel}

\usepackage[
    top=2.5cm,
    bottom=2.5cm,
    left=3cm,
    right=2.5cm,
    headheight=15pt
]{geometry}

\usepackage{amsmath}
\usepackage{amssymb}
\usepackage{mathtools}

\usepackage{graphicx}
\usepackage{float}
\usepackage{caption}
\usepackage{subcaption}

\usepackage{booktabs}
\usepackage{array}
\usepackage{longtable}
\usepackage{tabularx}

\usepackage{enumitem}

\usepackage{xcolor}
\usepackage[most]{tcolorbox}

\usepackage{fancyhdr}
\usepackage{ragged2e}

\usepackage[
    colorlinks=true,
    linkcolor=blue!50!black,
    citecolor=green!40!black,
    urlcolor=blue!60!black,
    pdftitle={\titulo},
    pdfauthor={\autor},
    pdfsubject={Apuntes universitarios},
    bookmarks=true
]{hyperref}

% ============================================================
% CONFIGURACIÓN
% ============================================================
\graphicspath{
    {./img/}
    {./graficos/}
    {./figuras/}
}

\setcounter{tocdepth}{2}

\setlength{\parindent}{0pt}
\setlength{\parskip}{0.5em}

\setlist[itemize]{
    topsep=0.3em,
    itemsep=0.2em
}

\setlist[enumerate]{
    topsep=0.3em,
    itemsep=0.2em
}

\clubpenalty=10000
\widowpenalty=10000

% ============================================================
% COMANDOS
% ============================================================
\newcommand{\abs}[1]{\left|#1\right|}
\newcommand{\norm}[1]{\left\lVert#1\right\rVert}
\newcommand{\vect}[1]{\mathbf{#1}}
\newcommand{\deriv}[2]{\frac{d#1}{d#2}}
\newcommand{\pderiv}[2]{\frac{\partial#1}{\partial#2}}

% Aviso para los textos normativos incluidos como referencia en el apunte
\newcommand{\complementario}{\par\smallskip{\small\itshape Texto de referencia complementario incluido en el apunte: no fue leído en clase. Conviene cotejarlo con el texto oficial actualizado (InfoLEG).}}

% Dimensiones de las tablas Cornell
\newcommand{\cornq}{0.27\textwidth}
\newcommand{\corna}{0.67\textwidth}

% ============================================================
% ENTORNOS / CAJAS
% ============================================================
\newtcolorbox{definition}[1]{
    enhanced,
    breakable,
    title={Definición: #1},
    fonttitle=\bfseries,
    colback=blue!3,
    colframe=blue!50!black,
    boxrule=0.8pt,
    arc=2mm
}

\newtcolorbox{important}{
    enhanced,
    breakable,
    title={Importante},
    fonttitle=\bfseries,
    colback=yellow!8,
    colframe=orange!70!black,
    boxrule=0.8pt,
    arc=2mm
}

\newtcolorbox{example}[1]{
    enhanced,
    breakable,
    title={Ejemplo: #1},
    fonttitle=\bfseries,
    colback=green!4,
    colframe=green!40!black,
    boxrule=0.8pt,
    arc=2mm
}

\newtcolorbox{formula}[1]{
    enhanced,
    breakable,
    title={Fórmula / Regla: #1},
    fonttitle=\bfseries,
    colback=purple!4,
    colframe=purple!50!black,
    boxrule=0.8pt,
    arc=2mm
}

\newtcolorbox{exercise}[1]{
    enhanced,
    breakable,
    title={Ejercicio: #1},
    fonttitle=\bfseries,
    colback=gray!5,
    colframe=gray!60!black,
    boxrule=0.8pt,
    arc=2mm
}

\newtcolorbox{note}{
    enhanced,
    breakable,
    title={Nota},
    fonttitle=\bfseries,
    colback=white,
    colframe=gray!50,
    boxrule=0.6pt,
    arc=2mm
}

% ============================================================
% CONFIGURACIÓN FINAL
% ============================================================
\pagestyle{fancy}
\fancyhf{}

\fancyhead[L]{\nouppercase{\leftmark}}
\fancyhead[R]{\materia}

\fancyfoot[C]{\thepage}

\renewcommand{\headrulewidth}{0.4pt}
\renewcommand{\footrulewidth}{0pt}

\fancypagestyle{plain}{
    \fancyhf{}
    \fancyfoot[C]{\thepage}
    \renewcommand{\headrulewidth}{0pt}
}

% ============================================================
% DOCUMENT
% ============================================================
\begin{document}

% ------------------------------------------------------------
% PORTADA
% ------------------------------------------------------------
\begin{titlepage}

\thispagestyle{empty}

\begin{center}

\vspace*{1cm}

{\large \textsc{\universidad}}

\vspace{0.2cm}

{\large \textsc{\facultad}}

\vspace{3cm}

{\Huge \textbf{\materia}}

\vspace{1cm}

{\LARGE \titulo}

\vspace{0.4cm}

{\normalsize \subtitulo}

\vspace{0.8cm}

\textit{Comprender el porqué de cada norma es el primer paso para convertir el estudio en conocimiento.}

\vspace{1.2cm}

\rule{0.70\textwidth}{0.5pt}

\vspace{0.5cm}

{\large Apuntes de estudio}

\vspace{0.5cm}

\rule{0.70\textwidth}{0.5pt}

\vfill

{\large \autor}

\vspace{0.3cm}

{\large \anio}

\end{center}

\vfill

\begin{flushleft}
\textit{Este es un modesto aporte a los estudiantes de ``\materia'' no pretende sustituir el material oficial de la materia.}
\end{flushleft}

\end{titlepage}

% ------------------------------------------------------------
% ÍNDICE
% ------------------------------------------------------------
\tableofcontents

% ------------------------------------------------------------
% CONTENIDO
% ------------------------------------------------------------

% ============================================================
\chapter[Repaso: Defensa del Consumidor]{Repaso: Defensa del Consumidor}
% ============================================================

\section{Encuadre: normas de regulación general y autoridades de aplicación}

El régimen de \textbf{defensa del consumidor} (Ley 24.240) y el de \textbf{defensa de la competencia} (Ley 27.442) constituyen dos grandes normas de regulación general de la actividad económica.

\begin{important}
Ambas leyes regulan la actividad económica \textbf{sin sesgo sectorial}; por ese motivo son normas de gran alcance en sus efectos. Cada una cuenta con su propia autoridad de aplicación.
\end{important}

En materia de defensa de la competencia, la ley previa preveía un tribunal estable que nunca llegó a consolidarse. Recientemente se ha constituido la autoridad prevista por la Ley 27.442, actualmente vigente, cuyos miembros son designados con acuerdo del Senado (el apunte consigna que, hasta el momento, fueron designados tres miembros). En materia de defensa del consumidor, la autoridad de aplicación a nivel nacional es la Secretaría de Comercio, a través de la Subsecretaría de Defensa del Consumidor.

\begin{table}[H]
\centering
\renewcommand{\arraystretch}{1.25}
\begin{tabularx}{\textwidth}{
    >{\raggedright\arraybackslash}p{0.2\textwidth}
    >{\raggedright\arraybackslash}X
    >{\raggedright\arraybackslash}X
}
\toprule
 & \textbf{Defensa del Consumidor} & \textbf{Defensa de la Competencia} \\
\midrule
\textbf{Norma} & Ley 24.240 & Ley 27.442 (actual) \\
\textbf{Autoridad de aplicación} & Secretaría de Comercio, Subsecretaría de Defensa del Consumidor (nivel nacional) & Tribunal / Autoridad de Defensa de la Competencia (miembros designados con acuerdo del Senado) \\
\textbf{Alcance} & Regulación general, sin sesgo sectorial & Regulación general, sin sesgo sectorial \\
\bottomrule
\end{tabularx}
\end{table}

\section{Estatus constitucional y concepto de consumidor}

\subsection{El estatus constitucional}

Con la reforma constitucional de 1994 se incorporó el artículo 42, que produjo la \textbf{constitucionalización de los derechos de usuarios y consumidores} y consolidó distintos aspectos de su tutela.

La defensa del consumidor, como verdadera política llevada adelante por el Estado, adquirió \textbf{estatus constitucional}. También lo tiene la defensa de la competencia, como se verá más adelante. La presencia de ambas en la Constitución indica que debe existir una verdadera política pública consecuente con la defensa de la parte débil, en general, de las relaciones económicas: los usuarios y consumidores.

% TODO: contenido incompleto o ambiguo en las notas originales
% (El apunte aclara que, en la primera mención, el docente dijo «la defensa de la competencia»; por el contexto se corrigió a «la defensa del consumidor».)

\begin{note}
\textbf{Constitución Nacional, art. 42 (primer y segundo párrafo).}
«Los consumidores y usuarios de bienes y servicios tienen derecho, en la relación de consumo, a la protección de su salud, seguridad e intereses económicos; a una información adecuada y veraz; a la libertad de elección, y a condiciones de trato equitativo y digno.»

«Las autoridades proveerán a la protección de esos derechos, a la educación para el consumo, a la defensa de la competencia contra toda forma de distorsión de los mercados, al control de los monopolios naturales y legales, al de la calidad y eficiencia de los servicios públicos, y a la constitución de asociaciones de consumidores y de usuarios.»
\complementario
\end{note}

\subsection{¿Quién es consumidor?}

\begin{definition}{Consumidor (según el docente)}
Son quienes se sirven de los bienes y servicios, contratándolos a título oneroso o gratuito, \textbf{para su consumo personal o de su grupo familiar}. No hay consumo cuando alguien contrata, por ejemplo, insumos o servicios incorporados a su actividad onerosa o comercial. Siempre se trata de lo que puede denominarse un consumo personal propio del grupo familiar.
\end{definition}

\section{La Ley 24.240 como ley de orden público}

\subsection{Carácter de orden público}

La Constitución otorgó a la materia un estatus que luego se plasma en la Ley 24.240 en un aspecto muy importante: se trata de una ley de \textbf{orden público}. Las normas a las que se asigna ese carácter deben prevalecer \textbf{aun cuando exista acuerdo en contrario entre las partes}, incluso un acuerdo voluntario.

La ley se destina a tutelar un desequilibrio en las relaciones económicas y a llevar a un pie de igualdad a dos partes que, por su naturaleza, probablemente llegan a este tipo de vínculo jurídico en un estado desigual: existe, en general, mayor necesidad e inexperiencia del consumidor. Esa es la razón por la que la ley recibe el carácter de orden público.

\begin{note}
\textbf{Ley 24.240, art. 65.} «Esta ley es de orden público, rige en todo el territorio nacional y entrará en vigencia a los noventa (90) días de su publicación en el Boletín Oficial.»
\complementario
\end{note}

\subsection{Consecuencia práctica: nulidad de cláusulas abusivas}

Sobre la base del orden público pueden establecerse nulidades en las cláusulas de los contratos de consumo. La doctrina y la jurisprudencia lo han desarrollado ampliamente, sobre todo en el fuero civil: aun cuando las partes hayan suscripto voluntariamente un contrato, si se encuentran cláusulas que a priori resultan abusivas, los jueces pueden considerarlas nulas y, por lo tanto, inaplicables. Esto vale especialmente para los contratos \textit{pro forma} o de adhesión.

\begin{example}{Cláusulas que el docente identifica como abusivas}
\begin{itemize}
    \item Las que \textbf{invierten la carga de la prueba} en desmedro del consumidor.
    \item Las que imponen al consumidor \textbf{renuncia de derechos}.
    \item Las que \textbf{dificultan el acceso a la justicia}.
    \item Las que \textbf{reducen irrazonablemente el ámbito de la responsabilidad} del proveedor del consumo.
\end{itemize}
\end{example}

\begin{note}
\textbf{Ley 24.240, art. 37 (incisos a, b y c).} «Sin perjuicio de la validez del contrato, se tendrán por no convenidas: a) Las cláusulas que desnaturalicen las obligaciones o limiten la responsabilidad por daños; b) Las cláusulas que importen renuncia o restricción de los derechos del consumidor o amplíen los derechos de la otra parte; c) Las cláusulas que contengan cualquier precepto que imponga la inversión de la carga de la prueba en perjuicio del consumidor.»
\complementario
\end{note}

\section{Cláusulas abusivas en la práctica: carteles y seguros}

\subsection{Un ejemplo «por el absurdo»: garajes, cines y shoppings}

Se trata de un ejemplo que hoy se encuentra bastante discutido en la doctrina, porque ha enfrentado a la jurisprudencia de la Corte y a la de la Cámara Nacional de Apelaciones en lo Civil.

\begin{example}{El cartel «no nos hacemos responsables»}
Un caso típico es el cartel que se encuentra en un garaje o en un cine, en el que el proveedor del servicio señala que el establecimiento no se hace responsable por el robo de objetos. Es una situación en la que se intenta desplazar la responsabilidad, limitándola irrazonablemente.

Si una persona deja su automóvil en un shopping y, al volver, se lo han robado, el proveedor no puede rehusar su responsabilidad por más que exhiba un cartel en el frente del estacionamiento, porque existe un \textbf{deber de guarda} y una relación de consumo que genera claramente responsabilidad. Además, ni siquiera se trata del ámbito contractual: es un simple aviso unilateral.
\end{example}

El propio docente aclara que el ejemplo se da casi por el absurdo.

\subsection{Un tema más complejo: el seguro obligatorio y su tope de cobertura}

En los contratos de seguro, por ejemplo el de responsabilidad civil, existe numerosa jurisprudencia según la cual, salvo que se trate de un seguro contra todo riesgo y de amplio espectro, el contrato tiene un \textbf{tope} de cobertura (por ejemplo, la garantía del vehículo hasta una suma determinada).

La Cámara Civil ha destacado que, por tratarse de un seguro de características muy generales y obligatorio para los conductores, debe primar el concepto de que es un \textbf{contrato efectivamente de consumo}: no lo utiliza un transportista para su trabajo, sino una persona para sus actividades cotidianas. En consecuencia, estableció que el tope de la garantía (la cobertura del seguro) no le es aplicable al damnificado, porque se trata de una restricción que limita la obligación de asegurar que tiene la aseguradora. Ante un siniestro, la aseguradora debería responder por lo que se establezca, salvo que la cobertura se corte por algún delito o situación como, por ejemplo, el consumo de alcohol al volante, y siempre que existan las condiciones generales para que el seguro preste su cobertura.

Las cámaras civiles han entendido, en su gran mayoría, que el tope es una restricción que no debe ser aplicable cuando hay un siniestro y se discute judicialmente. Se trata de una aplicación de la Ley 24.240, que considera el tope como \textbf{cláusula abusiva}.

\subsection{La respuesta de la Corte: norma especial sobre norma general}

La complejidad radica en que subyace a las obligaciones toda una ecuación económica: no se paga lo mismo por la prima de un seguro contra todo riesgo que por la de un contrato de seguro de menor porte.

La Corte ha dicho que, sin perjuicio del concepto general, debe primar la norma especial. En numerosos fallos hizo prevalecer la ley del seguro como ley específica, desplazó el criterio de la Cámara Civil y desautorizó muchas de sus sentencias. Su razonamiento es que no puede pretenderse una responsabilidad absoluta cuando el consumidor, al contratar, conoce de antemano los límites y paga una prima en retribución de la cobertura que posee. Debe primar, entonces, la \textbf{norma especial por sobre la norma general}. Queda por verse, con la actual composición de la Corte, qué postura se adopta; por ahora existe una controversia latente entre dos posturas.

\begin{table}[H]
\centering
\renewcommand{\arraystretch}{1.25}
\begin{tabularx}{\textwidth}{
    >{\raggedright\arraybackslash}p{0.2\textwidth}
    >{\raggedright\arraybackslash}X
    >{\raggedright\arraybackslash}X
}
\toprule
 & \textbf{Cámara Civil} & \textbf{Corte Suprema} \\
\midrule
\textbf{Enfoque} & Postura garantista, con enorme prevalencia y proyección de la ley de defensa del consumidor & Postura más realista, que considera racionalmente la ecuación económica que subyace al contrato \\
\textbf{Naturaleza del seguro} & Es sobre todo un contrato de consumo (seguro obligatorio, de uso cotidiano) & Es un contrato de seguro: rige la ley específica del seguro \\
\textbf{Tope de cobertura} & No es oponible al damnificado: es una cláusula abusiva & Es válido: el consumidor lo conoce de antemano y paga una prima proporcional a la cobertura \\
\textbf{Regla que invoca} & Aplicación de la Ley 24.240 & La norma especial prevalece sobre la norma general \\
\textbf{Estado actual} & Criterio mayoritario en las cámaras civiles & Ha desautorizado numerosas sentencias; falta ver qué postura adopta la actual composición \\
\bottomrule
\end{tabularx}
\end{table}

Quien paga una cobertura menor no puede pretender la misma cobertura que quien paga una mayor. Por otro lado, si se aplica la ley (en un artículo que siempre se destaca en esta controversia), la compañía de seguros deberá responder \textbf{en la medida del seguro} contratado. Existe además un vínculo contractual muy concreto y la intervención de la Superintendencia de Seguros como autoridad de control. Por lo tanto, es discutible entender si hay o no una suerte de aprovechamiento del consumidor inexperto.

\begin{note}
El apunte indica que la expresión «en la medida del seguro» corresponde al régimen de la Ley de Seguros (Ley 17.418, art. 118); el docente no consignó el número del artículo.
\end{note}

El ejemplo permite poner en perspectiva que subsisten controversias: una postura tiene un ojo mucho más puesto en la realidad económica y otra sostiene la prevalencia de la ley de defensa del consumidor. Lo que importa para el tema es que la ley de defensa del consumidor se mantiene más allá de lo acordado por las partes; por eso tiene carácter de orden público.

\section{Los grandes conceptos de la Ley 24.240}

\subsection{Información veraz, cierta, concreta y oportuna}

Un concepto clave de la ley es el derecho del consumidor a una información \textbf{veraz, concreta, cierta y oportuna}. El proveedor del bien de consumo debe informar sobre los riesgos, las características y el alcance. No puede haber información que lleve a confusión y, mucho menos, que pueda perjudicar; así ocurre, por ejemplo, con la información sobre bienes alimenticios y fármacos, temas de enorme sensibilidad donde está en juego un bien jurídico como la salud. La información cierta, veraz, concreta y accesible es un aspecto clave de la ley y debe primar.

\begin{note}
\textbf{Ley 24.240, art. 4 (primer párrafo).} «El proveedor está obligado a suministrar al consumidor en forma cierta, clara y detallada todo lo relacionado con las características esenciales de los bienes y servicios que provee, y las condiciones de su comercialización. La información debe ser siempre gratuita para el consumidor y proporcionada con la claridad necesaria que permita su comprensión.»
\complementario
\end{note}

\subsection{Trato digno y no discriminatorio}

La ley busca evitar los tratos discriminatorios y denigrantes y garantizar el \textbf{trato digno} del consumidor. Aunque es un concepto indeterminado, a partir de él pueden dilucidarse muchas controversias y constituye un estándar que fija obligaciones para el proveedor.

\begin{note}
\textbf{Ley 24.240, art. 8 bis (primeros párrafos).} «Los proveedores deberán garantizar condiciones de atención y trato digno y equitativo a los consumidores y usuarios. Deberán abstenerse de desplegar conductas que coloquen a los consumidores en situaciones vergonzantes, vejatorias o intimidatorias.»
\complementario
\end{note}

\subsection{Responsabilidad y solidaridad a lo largo de la cadena}

Otro aspecto clave es la responsabilidad y la \textbf{solidaridad} que presume y prevé la ley: alcanza a toda la cadena, desde la producción hasta la última etapa, incluso el transportista, aunque la responsabilidad de este es más acotada. La solidaridad tiende a lograr una fácil demandabilidad por parte del eventual consumidor damnificado. Su finalidad es que el consumidor tenga a quién demandar y que no exista una suerte de «pase de la pelota» para evitar responsabilidades ni conductas dilatorias en los procesos. El consumidor puede demandar a cualquier integrante de la cadena.

\begin{note}
\textbf{Ley 24.240, art. 40 (primer párrafo).} «Si el daño al consumidor resulta del vicio o riesgo de la cosa o de la prestación del servicio, responderán el productor, el fabricante, el importador, el distribuidor, el proveedor, el vendedor y quien haya puesto su marca en la cosa o servicio. El transportista responderá por los daños ocasionados a la cosa con motivo o en ocasión del servicio. La responsabilidad es solidaria, sin perjuicio de las acciones de repetición que correspondan. Sólo se liberará total o parcialmente quien demuestre que la causa del daño le ha sido ajena.»
\complementario
\end{note}

\subsection{La publicidad como parte de la oferta}

También es muy importante la información en el marco de la \textbf{publicidad}. Según la norma, lo que se enuncia en la publicidad de un producto forma parte, prácticamente, del prospecto. No puede existir información confusa, falaz, poco veraz o inexacta: todos son aspectos que hacen a la tutela de los derechos del consumidor.

\begin{note}
\textbf{Ley 24.240, art. 8.} «Las precisiones formuladas en la publicidad o en anuncios, prospectos, circulares u otros medios de difusión obligan al oferente y se tienen por incluidas en el contrato con el consumidor.»
\complementario
\end{note}

\section{Autoridad de aplicación, daño punitivo y acciones de clase}

\subsection{La autoridad de aplicación y el daño punitivo}

Entre los aspectos clave de la ley se cuentan la autoridad de aplicación, con la posibilidad de dirimir conflictos e imponer sanciones, y la posibilidad de imponer sanciones de carácter preventivo o disciplinador del mercado. Se puede tener en cuenta la posición del proveedor en el mercado (por ejemplo, un proveedor importante, con una gran porción de ese mercado) y la reiteración de conductas violatorias del consumo o de la ley de defensa del consumidor. En esos casos la ley prevé el \textbf{daño punitivo}, destinado a desalentar las prácticas que generan conflicto y perjuicios al consumidor. Es un aspecto novedoso que ya tiene mucho recorrido.

\begin{note}
\textbf{Ley 24.240, art. 52 bis (daño punitivo).} «Al proveedor que no cumpla sus obligaciones legales o contractuales con el consumidor, a instancia del damnificado, el juez podrá aplicar una multa civil a favor del consumidor, la que se graduará en función de la gravedad del hecho y demás circunstancias del caso, independientemente de otras indemnizaciones que correspondan. Cuando más de un proveedor sea responsable del incumplimiento responderán todos solidariamente ante el consumidor, sin perjuicio de las acciones de regreso que les correspondan. La multa civil que se imponga no podrá superar el máximo de la sanción de multa prevista en el artículo 47, inciso b) de esta ley.»
\complementario
\end{note}

\subsection{El registro de asociaciones de defensa del consumidor}

La Secretaría de Comercio es organismo de aplicación también del registro de acciones de clase. La Subsecretaría de Defensa del Consumidor es la autoridad del \textbf{registro de asociaciones de defensa del consumidor}. Esto es importante porque las asociaciones registradas son las que tienen legitimación para accionar en defensa de estos derechos.

\subsection{La legitimación activa y el fallo «Halabi»}

La ley, junto con el fallo «Halabi», estableció, sobre todo a partir de la reforma constitucional, una ampliación de la \textbf{legitimación activa}. El docente menciona además, como referencia, el tema de los excombatientes de Malvinas y el fallo de la Corte, muy discutido.

\begin{definition}{Legitimación activa}
Es un presupuesto elemental para una acción judicial, para que haya un verdadero proceso judicial: tiene que haber alguien con aptitud legal para demandar, a partir de un interés o un derecho reconocible en el universo jurídico y que pueda ser invocado.
\end{definition}

La Constitución, desde su ampliación, y la normativa moderna que amplió la legitimación activa establecen que las asociaciones de defensa del consumidor pueden accionar en defensa de los consumidores. En las denominadas \textbf{acciones de clase} pueden accionar representando a la clase, siempre que se trate de \textbf{derechos individuales homogéneos}: son homogéneos porque tienen en común un mismo hecho generador del daño.

\begin{note}
\textbf{Constitución Nacional, art. 43 (segundo párrafo).} «Podrán interponer esta acción contra cualquier forma de discriminación y en lo relativo a los derechos que protegen al ambiente, a la competencia, al usuario y al consumidor, así como a los derechos de incidencia colectiva en general, el afectado, el defensor del pueblo y las asociaciones que propendan a esos fines, registradas conforme a la ley, la que determinará los requisitos y formas de su organización.»
\complementario
\end{note}

\begin{example}{El banco y la promoción engañosa}
Un gran banco lanza una promoción que termina engañando y perjudicando a los consumidores: todos los clientes del banco podrían accionar. Quizás individualmente el reclamo no sea de gran entidad, pero al sumar todos los reclamos hay un incentivo para que una asociación lo promueva. La ley persigue una herramienta jurídica para combatir a las grandes empresas y corporaciones que a veces se aprovechan de consumidores inexpertos, atomizados y no organizados para reclamar por sus derechos.
\end{example}

\subsection{Registro nacional y registros locales}

En el marco de la Subsecretaría de Defensa del Consumidor existe el registro de asociaciones de defensa del consumidor. Las asociaciones inscriptas pueden demandar en procesos colectivos, en acciones de clase y en defensa de derechos del consumidor a nivel nacional. También existen registros locales: para demandar ante el fuero nacional o federal se requiere la habilitación en el registro nacional; para demandar solo en el ámbito de la Ciudad de Buenos Aires puede alcanzar la habilitación en el registro de asociaciones de defensa del consumidor de la Ciudad.

La autoridad de aplicación de la ley resuelve conflictos entre las partes, puede imponer sanciones y es autoridad de aplicación del registro. Ninguna asociación de defensa del consumidor podría intervenir en acciones de clase si no está incluida, a nivel nacional, en ese registro.

\subsection{Anticipo de temas posteriores}

El docente anticipa que, además de la defensa de la competencia, se abordará el decreto-ley, convalidado por el Congreso, que establece las sanciones para los actos de \textbf{lealtad comercial} y que ha sufrido un fallo adverso.

\section{Cuadro Cornell del capítulo}

\begin{longtable}{@{}>{\raggedright\arraybackslash}p{\cornq}!{\vrule width 0.4pt}>{\raggedright\arraybackslash}p{\corna}@{}}
\toprule
\textbf{Preguntas / palabras clave} & \textbf{Notas / respuestas} \\
\midrule
\endfirsthead
\toprule
\textbf{Preguntas / palabras clave} & \textbf{Notas / respuestas} \\
\midrule
\endhead
\midrule
\multicolumn{2}{r}{\small\itshape continúa en la página siguiente} \\
\endfoot
\bottomrule
\endlastfoot

\textbf{¿Por qué la defensa del consumidor tiene estatus constitucional y quién es consumidor?} &
Porque el art. 42 de la Constitución (reforma de 1994) la incorporó como política pública de protección de la parte débil de las relaciones económicas. Es consumidor quien contrata bienes o servicios, incluso gratuitamente, \textbf{para su consumo personal o de su grupo familiar}; no lo es quien contrata insumos para su actividad comercial. \\
\addlinespace[0.6em]

\textbf{¿Qué significa que la Ley 24.240 sea de orden público?} &
Que prevalece \textbf{aun cuando haya acuerdo en contrario} entre las partes. Permite declarar nulas las cláusulas abusivas (inversión de la carga de la prueba, renuncia de derechos, obstáculos al acceso a la justicia, limitación irrazonable de la responsabilidad), sobre todo en contratos de adhesión. \\
\addlinespace[0.6em]

\textbf{¿Qué valor tiene el cartel «no nos hacemos responsables»?} &
Ninguno frente a la relación de consumo: es un simple aviso unilateral que intenta limitar irrazonablemente la responsabilidad. En un garaje o en un estacionamiento de shopping hay un \textbf{deber de guarda}. \\
\addlinespace[0.6em]

\textbf{¿Cuál es la controversia sobre el tope de cobertura del seguro?} &
La Cámara Civil lo considera un contrato de consumo y el tope, una cláusula abusiva inaplicable. La Corte hace prevalecer la \textbf{norma especial (ley de seguros)} sobre la general: el consumidor conoce los límites y paga una prima proporcional; la aseguradora responde «en la medida del seguro». \\
\addlinespace[0.6em]

\textbf{¿Qué exige la ley en información y publicidad?} &
Información \textbf{cierta, veraz, concreta, accesible y oportuna}, sobre todo en alimentos y fármacos (salud). Lo que dice la publicidad integra la oferta y obliga al proveedor; no puede ser confusa, falaz ni inexacta. \\
\addlinespace[0.6em]

\textbf{¿Qué son el trato digno y la solidaridad?} &
El trato digno prohíbe los tratos discriminatorios y denigrantes y fija obligaciones al proveedor. La responsabilidad es \textbf{solidaria en toda la cadena} (del productor al vendedor, con el transportista de modo más acotado), para que el consumidor tenga a quién demandar y se evite el «pase de la pelota». \\
\addlinespace[0.6em]

\textbf{¿Para qué sirve el daño punitivo?} &
Para \textbf{desalentar} las prácticas que perjudican al consumidor, considerando la posición del proveedor en el mercado y la reiteración de conductas violatorias. \\
\addlinespace[0.6em]

\textbf{¿Quién puede promover acciones de clase y por qué?} &
Las \textbf{asociaciones de defensa del consumidor inscriptas} en el registro (nacional, o local para la Ciudad). Tienen legitimación activa ampliada (Constitución y fallo «Halabi») cuando hay \textbf{derechos individuales homogéneos}, es decir, un mismo hecho generador del daño (por ejemplo, la promoción engañosa de un banco). \\

\midrule
\multicolumn{2}{@{}p{\dimexpr\cornq+\corna+2\tabcolsep\relax}@{}}{%
\textbf{Resumen:} La defensa del consumidor tiene estatus constitucional (art. 42) y se plasma en la Ley 24.240, de orden público, que prevalece sobre los pactos y permite anular cláusulas abusivas. Sus grandes conceptos son la información veraz, el trato digno, la solidaridad en la cadena de consumo y la publicidad como parte de la oferta; la ley prevé además el daño punitivo y las acciones de clase promovidas por asociaciones registradas.} \\
\end{longtable}

% ============================================================
\chapter[Defensa de la Competencia: bases]{Defensa de la Competencia: bases}
% ============================================================

\section{Concepto y finalidad de la defensa de la competencia}

\subsection{Una competencia genuina y real}

La defensa de la competencia busca que haya, en definitiva, una \textbf{competencia genuina, una competencia real}. El docente precisa que esta expresión resulta más ajustada que la de «competencia justa entre empresas».

\subsection{Evitar las distorsiones del mercado}

Es una ley que procura evitar distorsiones del mercado. El mercado cumple un determinado rol como organizador de la oferta y la demanda; pero, más allá de algunas posiciones fanáticas, a grandes rasgos se reconoce que el mercado tiene disfuncionalidades y puede presentar distorsiones, entre ellas el \textbf{abuso} de aquellos agentes dominantes de la economía que tienen una posición fuerte y prevalente en el mercado.

La moderna teoría económica, aun bajo un paradigma liberal o de libre mercado, establece que esta ley es una herramienta necesaria, destinada a actuar cuando existen excesivas concentraciones del mercado que pueden perjudicar a consumidores y usuarios, es decir, cuando se produce lo que en la jerga se denomina \textbf{apropiación de la renta del consumidor}: una posición dominante puede generar mayores costos artificialmente.

\section{Las dos vertientes de la ley: preventiva y sancionatoria}

La ley debe distinguirse en dos aspectos clave: tiene una vertiente de carácter \textbf{preventivo} y otra de carácter \textbf{sancionatorio}.

La ley intenta evitar que haya una excesiva concentración en el mercado, porque de esas concentraciones pueden surgir conductas disvaliosas y distorsivas que, en esencia, perjudican a los consumidores. Por eso actúa con carácter preventivo, o bien actúa cuando encuentra una violación a sus preceptos, es decir, una conducta anticompetitiva. La vertiente preventiva es la que menos se conoce, pero también es muy importante.

\begin{table}[H]
\centering
\renewcommand{\arraystretch}{1.25}
\begin{tabularx}{\textwidth}{
    >{\raggedright\arraybackslash}p{0.2\textwidth}
    >{\raggedright\arraybackslash}X
    >{\raggedright\arraybackslash}X
}
\toprule
 & \textbf{Vertiente preventiva} & \textbf{Vertiente sancionatoria} \\
\midrule
\textbf{Qué busca} & Evitar lo que todavía no se ha configurado: la excesiva concentración del mercado & Sancionar conductas anticompetitivas ya cometidas \\
\textbf{Cuándo actúa} & Ante una concentración económica (por ejemplo, una empresa que compra a otra de su competencia) & Cuando hay una denuncia o una investigación y se demuestra la infracción \\
\textbf{Decisión posible} & Aprobar, condicionar o rechazar la concentración & Sancionar (multa) o no, tras un procedimiento administrativo \\
\textbf{¿Se sanciona a alguien?} & No: en esta vertiente no se sanciona a nadie & Sí, si se demuestra la conducta \\
\textbf{Imagen del docente} & Control previo de las concentraciones & «Una suerte de juzgado penal de instrucción» \\
\bottomrule
\end{tabularx}
\end{table}

\section{La autoridad de aplicación: composición e independencia}

\subsection{Cinco miembros con acuerdo del Senado}

La ley prevé una autoridad de aplicación: el Tribunal de Defensa de la Competencia o, mejor dicho, la \textbf{Autoridad de Defensa de la Competencia}. Está compuesta por cinco miembros que actúan en la órbita del Poder Ejecutivo pero que cuentan con acuerdo del Senado; por lo tanto, una vez nombrados adquieren inamovilidad en el cargo, salvo mala conducta, y no pueden ser desplazados arbitrariamente, lo que les da una condición de independencia.

\begin{important}
\textbf{Composición de la Autoridad de Defensa de la Competencia}
\begin{itemize}
    \item Cinco miembros: un presidente y cuatro vocales.
    \item Dos de ellos con especialidad de abogado y dos en ciencias económicas.
    \item Elegidos luego de un concurso ante autoridades académicas, a propuesta del Ejecutivo pero siempre con la convalidación del Congreso: en definitiva, como si fuera un magistrado, con el acuerdo del Senado.
    \item Inamovilidad en el cargo, salvo mala conducta: no pueden ser desplazados arbitrariamente.
\end{itemize}
\end{important}

\subsection{Un cambio profundo respecto de la ley anterior}

La existencia de una autoridad con acuerdo del Senado es un viejo anhelo desde la ley anterior; recién hace poco fueron designados los primeros tres miembros de la actual autoridad (de las cinco previstas, al menos tres están designadas). Hasta ese momento no había autoridades administrativas antitrust con acuerdo del Senado. Las decisiones de autorizar una concentración o de multar, que son las dos grandes ramas de la ley, eran tomadas por el secretario de Comercio, y la autoridad de la competencia tenía solo un carácter [consultivo, no vinculante]: emitía un dictamen o una recomendación.

% TODO: contenido incompleto o ambiguo en las notas originales
% (El apunte aclara que el audio dice «si bien vinculante», pero la frase siguiente describe un rol consultivo; se consignó entre corchetes la lectura coherente con el contexto.)

\section{Cuadro Cornell del capítulo}

\begin{longtable}{@{}>{\raggedright\arraybackslash}p{\cornq}!{\vrule width 0.4pt}>{\raggedright\arraybackslash}p{\corna}@{}}
\toprule
\textbf{Preguntas / palabras clave} & \textbf{Notas / respuestas} \\
\midrule
\endfirsthead
\toprule
\textbf{Preguntas / palabras clave} & \textbf{Notas / respuestas} \\
\midrule
\endhead
\midrule
\multicolumn{2}{r}{\small\itshape continúa en la página siguiente} \\
\endfoot
\bottomrule
\endlastfoot

\textbf{¿Para qué sirve la defensa de la competencia?} &
Para asegurar una \textbf{competencia genuina y real}, evitando las distorsiones y el abuso de los agentes dominantes, que pueden encarecer artificialmente y apropiarse de la renta del consumidor. \\
\addlinespace[0.6em]

\textbf{¿Cuáles son las dos vertientes de la ley?} &
\textbf{Preventiva:} control de concentraciones antes de que el daño se configure (no se sanciona a nadie). \textbf{Sancionatoria:} investigación y castigo de conductas anticompetitivas, con un procedimiento similar al de un juzgado de instrucción. \\
\addlinespace[0.6em]

\textbf{¿Cómo se compone la Autoridad y por qué es independiente?} &
Cinco miembros (presidente y cuatro vocales; dos abogados y dos en ciencias económicas), elegidos por concurso, propuestos por el Ejecutivo y con \textbf{acuerdo del Senado}, con inamovilidad salvo mala conducta. Antes decidía el secretario de Comercio y la comisión solo dictaminaba. \\

\midrule
\multicolumn{2}{@{}p{\dimexpr\cornq+\corna+2\tabcolsep\relax}@{}}{%
\textbf{Resumen:} La defensa de la competencia procura una competencia real y evitar las distorsiones del mercado y el abuso de los agentes dominantes. La ley tiene una vertiente preventiva (control de concentraciones) y otra sancionatoria (conductas anticompetitivas). Su autoridad de aplicación, de cinco miembros con acuerdo del Senado e inamovilidad, representa un cambio profundo respecto del régimen anterior.} \\
\end{longtable}

% ============================================================
\chapter[Vertiente preventiva: control de concentraciones]{Vertiente preventiva: el control de concentraciones}
% ============================================================

\section{Concentraciones económicas y umbral de intervención}

\subsection{Qué es una concentración y cuándo interviene la Autoridad}

La vertiente preventiva busca evitar lo que todavía no se ha configurado. Cuando hay concentraciones económicas (es decir, cuando el control de una unidad de negocios o empresa se concentra en una misma figura), la Autoridad de Defensa de la Competencia puede tener que intervenir si las concentraciones superan las [cien millones de (100.000.000)] \textbf{unidades móviles}, una unidad de valor que se va actualizando. La Autoridad debe aprobar o desestimar la concentración económica.

Esto ocurre generalmente cuando una empresa compra a otra de su competencia. Si se trata de un volumen de negocios importante y económicamente relevante, la Autoridad debe resolver si esa adquisición, que genera mayor concentración en un mercado, es válida o si puede haber peligros para la competencia por una excesiva concentración.

Esta es la competencia de carácter preventivo de la autoridad de aplicación: una empresa que compra a otra, o un grupo económico que se fusiona y cambia de manos a sus dueños, puede generar una excesiva concentración dentro de un mercado relevante. Para ello es muy importante, en primer lugar, determinar cuál es el \textbf{mercado relevante}.

\begin{note}
\textbf{Ley 27.442, art. 7 (concentración económica) y art. 9 (umbral).}
Art. 7: «A los efectos de esta ley se entiende por concentración económica la toma de control de una o varias empresas, a través de la realización de los siguientes actos: a) La fusión entre empresas; b) La transferencia de fondos de comercio; [\dots]».
Art. 9 (primer párrafo, en lo pertinente): «Los actos indicados en el artículo 7° de la presente ley, cuando la suma del volumen de negocio total del conjunto de empresas afectadas supere en el país la suma equivalente a cien millones (100.000.000) de unidades móviles, deberán ser [notificados para su examen por el Tribunal de Defensa de la Competencia]».
\complementario{} El texto entre corchetes completa la cita de memoria de la norma; verificar la redacción oficial.
\end{note}

% TODO: contenido incompleto o ambiguo en las notas originales
\begin{note}
El apunte señala que, en el audio, el docente dice «100.000 unidades móviles», y que el umbral del art. 9 de la Ley 27.442 es de \textbf{cien millones (100.000.000)} de unidades móviles; se consignó entre corchetes el valor legal. El valor de la unidad móvil se actualiza periódicamente.
\end{note}

\section{Caso Bimbo--Fargo: aprobar, condicionar o rechazar}

\begin{example}{Bimbo compró Fargo}
Un caso de hace muchos años: Bimbo, gran multinacional de origen mexicano con alcance global, compró Fargo en la Argentina. Fargo era una empresa muy fuerte como marca y, aunque cambió de manos, mantuvo su marca. El mercado en cuestión es el de los panificados: el pan lactal envasado que se compra en el supermercado y que tiene un plazo de duración algo más amplio que el pan de panadería.

En ese producto, Bimbo, que ya tenía una participación en el mercado con sus importaciones, al comprar a Fargo (el gran actor del mercado) más algunas marcas secundarias, quedó prácticamente como dueña de la gran mayoría del mercado. Esa concentración económica debió someterse a una resolución de la autoridad de competencia (que entonces se llamaba «comisión»; hoy es «autoridad»). Se emitió un dictamen que autorizó la operación, pero con condicionamientos: vender algunas segundas marcas, establecer algunas restricciones de comercialización, entre otros.
\end{example}

Ante una concentración, la Autoridad puede aprobarla, condicionar la adquisición o rechazarla. El rechazo está reservado para casos extremos de una gran concentración económica.

\begin{table}[H]
\centering
\renewcommand{\arraystretch}{1.25}
\begin{tabularx}{\textwidth}{
    >{\raggedright\arraybackslash}p{0.18\textwidth}
    >{\raggedright\arraybackslash}X
    >{\raggedright\arraybackslash}X
}
\toprule
\textbf{Decisión de la Autoridad} & \textbf{Qué significa} & \textbf{Cuándo se usa (según el docente)} \\
\midrule
\textbf{Aprobar} & «Marche, todo bien, no hay problema» & Cuando la operación no preocupa para la competencia \\
\textbf{Condicionar} & Se aprueba, pero bajo ciertas reglas & Caso Bimbo--Fargo: vender segundas marcas, restricciones de comercialización; caso Brahma--Quilmes \\
\textbf{Rechazar} & Se impide la adquisición & Casos extremos de una gran concentración económica \\
\bottomrule
\end{tabularx}
\end{table}

\section{La posición dominante}

La finalidad de la ley es evitar el \textbf{abuso de la posición dominante}.

\begin{definition}{Posición dominante (según el docente)}
Una empresa tiene posición dominante cuando, por su grado de participación en el mercado, puede prácticamente fijar las condiciones de venta: posee una porción del mercado tan grande que casi no está expuesta a la competencia y, por lo tanto, puede poner el precio que quiera, al ser prácticamente la única empresa que comercializa y tiene trascendencia en ese mercado.
\end{definition}

En esta vertiente preventiva no se sanciona a nadie: lo que se hace es aprobar o rechazar una concentración, porque se quiere evitar de antemano una excesiva concentración.

\section{Cada mercado se analiza por separado: barreras de entrada y salida}

\subsection{El caso Cablevisión--Multicanal}

Cuando Cablevisión compró a Multicanal se hablaba de una concentración económica que ciertamente existía. La defensa de la competencia tiene muchísimos aspectos; lo que importa comprender es que \textbf{cada mercado se investiga y se analiza con sus particularidades}.

\subsection{Las barreras de entrada y salida}

Al analizar si existe o no posición dominante, un punto clave son las \textbf{barreras de entrada y salida} del mercado. Hay mercados más restrictivos que otros. En el caso de los medicamentos, por ejemplo, no es posible instalar una fábrica y salir a vender: se requiere un centro habilitado por la ANMAT, patentes y una serie de condiciones. Las eléctricas son todavía más reguladas, porque son servicio público; el docente anticipa que se hablará de qué significa el servicio público, probablemente en la próxima clase. El servicio público es el tipo de mercado más regulado que existe, porque en él se fija la tarifa. Hay también mercados que, sin llegar a ser servicio público ni tener tarifa fija de antemano, son muy regulados y tienen muchas barreras de entrada y salida.

Por eso, una empresa puede tener una participación muy alta en un mercado en el que prácticamente cualquiera puede entrar con algo de dinero, y preocupar menos. En cambio, en otros mercados la participación puede ser más baja pero, por la cantidad de requisitos, preocupar más; es el caso de las aerolíneas, donde hay que tener una ruta y ganar una concesión, y el porte económico es mucho más alto. Cada mercado se analiza individualmente.

\begin{table}[H]
\centering
\renewcommand{\arraystretch}{1.25}
\begin{tabularx}{\textwidth}{
    >{\raggedright\arraybackslash}p{0.2\textwidth}
    >{\raggedright\arraybackslash}X
    >{\raggedright\arraybackslash}X
}
\toprule
\textbf{Tipo de mercado} & \textbf{Ejemplos del docente} & \textbf{Cómo pesa la participación} \\
\midrule
\textbf{Con barreras de entrada y salida bajas} & Cervezas (prácticamente pueden entrar y salir productores sin grandes requisitos) & Una participación muy alta preocupa menos: cualquiera puede entrar «con un poco de plata» \\
\textbf{Con barreras altas (muy regulados)} & Medicamentos (habilitación de la ANMAT, patentes); aerolíneas (rutas, concesiones) & Aun con participación más baja puede haber preocupación: son muchos los requisitos para entrar \\
\textbf{Servicio público} & Distribución eléctrica & Es el tipo de mercado más regulado que existe: se fija la tarifa \\
\bottomrule
\end{tabularx}
\end{table}

\section{Brahma--Quilmes y el mercado geográfico}

\subsection{Las tres respuestas posibles de la Autoridad}

La Autoridad de Defensa de la Competencia puede autorizar (considerando que no hay problema en ese mercado), condicionar («apruebo, pero bajo ciertas reglas») o directamente rechazar. El rechazo procede cuando se trata de una situación muy grave de concentración.

\subsection{Brahma compra Quilmes}

\begin{example}{Brahma--Quilmes}
Brahma, una multinacional que el docente cree de origen europeo (expresa dudas al respecto: «si no me equivoco») y gran marca de cervezas, compró Quilmes, con lo cual alcanzó una participación enorme en el mercado. Si bien las cervezas son productos relativamente desregulados (pueden entrar y salir productores sin grandes requisitos, más allá de las condiciones sanitarias normales de una bebida), había una gran concentración. Se autorizó la compra, pero con algunos condicionamientos.
\end{example}

\begin{note}
El apunte aclara que el grupo que adquirió Quilmes era AmBev, de origen brasileño. La referencia a «Kilmes» del audio se normalizó como «Quilmes».
\end{note}

\subsection{El mercado también se analiza geográficamente: las cadenas de cines}

\begin{example}{Una cadena de cines compra a otra}
El docente recuerda que, hace un tiempo, una cadena grande de cines compró a otra (no recuerda cuáles). También allí hubo que analizar el mercado geográficamente: en la Ciudad de Buenos Aires hay muchos shoppings y muchos cines, comunes y de shopping; pero en una provincia importante, como Santa Fe, probablemente no hay tantos cines. El análisis debe hacerse en el mercado relevante, con las condiciones de ese mercado y también geográficamente.
\end{example}

\section{Integración vertical: aguas arriba y aguas abajo}

Otro punto que no debe pasarse por alto en las concentraciones es la compra de una empresa que quizás no está en el mismo mercado, sino \textbf{aguas abajo o aguas arriba}.

\begin{example}{Coca-Cola compra la embotelladora}
Si Coca-Cola comprara la única gran embotelladora del país (el docente aclara que exagera), sería preocupante para la competencia: ¿qué ocurriría si Coca-Cola dejara de hacer las botellas a todos los competidores para quedarse con el mercado?
\end{example}

En relación con esto, se recuerda que en el caso de Quilmes hubo un problema porque tenía exclusividad con bares. Puede haber algunas promociones, pero no puede haber lo que se denomina una \textbf{exclusión total} del mercado.

\subsection{Síntesis de la vertiente preventiva}

La Autoridad de Defensa de la Competencia, con cinco miembros que dependen del Poder Ejecutivo pero tienen estabilidad en el cargo, puede aprobar, condicionar o rechazar una concentración económica, cuidando de que no se configure una posición dominante.

\begin{important}
\textbf{Posición dominante vs. abuso de posición dominante.} La posición dominante preocupa cuando hay un abuso. La posición dominante \textit{per se} no es deseada ni es lo ideal, pero una empresa puede tener todo el mercado y conducirse sin agraviar a los consumidores. Lo que constituye una infracción es abusar de esa posición: excluir competidores y tomar maniobras para quedarse con todo el mercado aprovechando la posición dominante.
\end{important}

Lo que se intenta evitar son las concentraciones excesivas; la otra gran misión de la ley es \textbf{sancionar conductas anticompetitivas}.

\section{Cuadro Cornell del capítulo}

\begin{longtable}{@{}>{\raggedright\arraybackslash}p{\cornq}!{\vrule width 0.4pt}>{\raggedright\arraybackslash}p{\corna}@{}}
\toprule
\textbf{Preguntas / palabras clave} & \textbf{Notas / respuestas} \\
\midrule
\endfirsthead
\toprule
\textbf{Preguntas / palabras clave} & \textbf{Notas / respuestas} \\
\midrule
\endhead
\midrule
\multicolumn{2}{r}{\small\itshape continúa en la página siguiente} \\
\endfoot
\bottomrule
\endlastfoot

\textbf{¿Cuándo interviene la Autoridad en una concentración y qué decide?} &
Cuando la operación supera el umbral de \textbf{cien millones de unidades móviles}. Puede \textbf{aprobar, condicionar o rechazar} (Bimbo--Fargo: condicionamientos; Brahma--Quilmes: autorizada con condicionamientos; el rechazo, para casos extremos). \\
\addlinespace[0.6em]

\textbf{¿Qué es la posición dominante y cuándo es infracción?} &
Tener una porción tan grande del mercado que casi no se está expuesto a la competencia y se pueden fijar las condiciones de venta. Por sí misma «no es lo ideal», pero \textbf{la infracción es el abuso}: excluir competidores para quedarse con todo el mercado. \\
\addlinespace[0.6em]

\textbf{¿Qué factores se analizan en cada mercado?} &
(1) El mercado relevante, incluido el ámbito \textbf{geográfico}; (2) las barreras de entrada y salida y el grado de regulación (medicamentos, eléctricas y aerolíneas frente a cervezas); (3) el número de competidores y su porcentaje; (4) la \textbf{sustitución} por otros productos. También se considera la integración \textbf{aguas arriba o abajo} (Coca-Cola y la embotelladora) y se rechaza la exclusión total. \\

\midrule
\multicolumn{2}{@{}p{\dimexpr\cornq+\corna+2\tabcolsep\relax}@{}}{%
\textbf{Resumen:} La vertiente preventiva controla las concentraciones económicas que superan el umbral de unidades móviles y no sanciona a nadie: la Autoridad aprueba, condiciona o rechaza. Cada mercado se analiza en sus particularidades (mercado relevante, ámbito geográfico, barreras de entrada y salida, integración vertical). La posición dominante preocupa cuando hay abuso.} \\
\end{longtable}

% ============================================================
\chapter[Vertiente sancionatoria: conductas anticompetitivas]{Vertiente sancionatoria: conductas anticompetitivas}
% ============================================================

\section{El caso de las cementeras}

\begin{example}{Cementeras: el primer gran fallo con multa}
Fue un fallo de hace muchísimos años y el \textbf{primer gran fallo en que se impuso una multa por la ley de defensa de la competencia}. A lo largo de un proceso (en el que se incautaron computadoras y hubo testigos) se demostró que las grandes empresas que proveían el cemento se reunían en la Cámara del Cemento Portland, en la Ciudad de Buenos Aires, y acordaban y se distribuían: acordaban el alza de precios (todas subían el precio al mismo tiempo para evitar competir y aumentar artificialmente sus ganancias) y se repartían las zonas geográficas para evitar competir.
\end{example}

Tras una investigación y un fallo que fue impugnado y tardó mucho en resolverse, terminó confirmándose una gran multa a las empresas.

\section{Ejemplos de conductas anticompetitivas}

\subsection{La venta atada}

En un supermercado, una cosa es una promoción y otra cosa es que no se pueda comprar un producto sin estar obligado a comprar otro. Es la denominada \textbf{venta atada}, ejemplo típico de conducta anticompetitiva.

\subsection{La cartelización: el caso de las prepagas}

\begin{example}{Aumento de las prepagas}
Un ejemplo típico de cartelización: nos ponemos de acuerdo y todos subimos los precios. Hace poco, al inicio de la gestión política vigente al momento de la clase, hubo un aumento de las prepagas que dio lugar al inicio de una investigación que, según el docente, todavía no había culminado. Lo que se sostuvo es que las prepagas se ponían de acuerdo para no generar competencia, pero subir todas el precio a un mismo nivel para ganar sin competir.
\end{example}

\subsection{Reparto de zonas, restricción de cantidades y bloqueo de competidores}

Es conducta anticompetitiva típica la coordinación o cartelización de precios y el \textbf{reparto de zonas} («en esta provincia vos no te metas, y yo no me meto en la otra, y cada uno en la provincia pone el precio que quiere»): no se genera competencia, y ello constituye un abuso y una violación de la ley. También lo son la restricción de cantidades y volúmenes para que el precio suba artificialmente en el mercado y la generación de ciertas prohibiciones a avances de empresas o de la competencia.

\begin{table}[H]
\centering
\renewcommand{\arraystretch}{1.25}
\begin{tabularx}{\textwidth}{
    >{\raggedright\arraybackslash}p{0.22\textwidth}
    >{\raggedright\arraybackslash}X
    >{\raggedright\arraybackslash}X
}
\toprule
\textbf{Conducta anticompetitiva} & \textbf{Ejemplo del docente} & \textbf{Efecto buscado} \\
\midrule
\textbf{Venta atada} & En un supermercado, no poder comprar un producto sin estar obligado a comprar otro (distinto de una simple promoción) & Atar la compra de un producto a la de otro \\
\textbf{Cartelización / coordinación de precios} & Cementeras; aumento de las prepagas & Subir precios todos a la vez para no competir \\
\textbf{Reparto de zonas geográficas} & «En esta provincia vos no te metas, y yo no me meto en la otra» & Cada uno fija el precio que quiere en su zona \\
\textbf{Restricción de cantidades y volúmenes} & Limitar la oferta & Que el precio suba artificialmente \\
\textbf{Bloqueo del avance de competidores} & Prohibiciones a empresas o a la competencia & Evitar que otros compitan \\
\bottomrule
\end{tabularx}
\end{table}

Todas estas acciones tienden a evitar un mercado donde los oferentes compitan de buena fe y razonablemente, lo que genera bienestar y mayor ventaja para el consumidor.

\section{Precios predatorios}

Aun cuando pueda parecer contradictorio, existe también una práctica anticompetitiva que consiste en ingresar a un precio demasiado barato: los \textbf{precios predatorios}. A propósito del comentario sobre las aerolíneas \textit{low cost}, el docente señala que es una hipótesis comentada durante mucho tiempo (que empresas europeas buscaban pisar la tarifa para hacer quebrar o desequilibrar a otras aerolíneas más grandes), aunque no sabe si está muy demostrada; algunos entienden que se trata, en realidad, de diferencias por flexibilidades, pero ha hecho mucho ruido en el mercado.

\begin{definition}{Precios predatorios}
\textbf{Conducta tolerada:} entrar al mercado con un precio muy bueno para darse a conocer. La literatura y la práctica de defensa de la competencia toleran esta actividad, porque generar alguna promoción por debajo del precio de producción para posicionarse viene bien para que nuevos actores ingresen y se hagan conocer.

\textbf{Conducta anticompetitiva:} aplicar precios por debajo de los propios costos solamente para hacer fundir al competidor, o mantener excesivamente en el tiempo un precio artificialmente bajo que no permita competir realmente a los adversarios del mercado. Es una conducta penada por la ley, aunque debe demostrarse.
\end{definition}

\begin{example}{El productor argentino y la multinacional}
Un productor argentino anda muy bien en un mercado. Llega una multinacional con mucho respaldo financiero que piensa: poniéndole el precio a mitad de costo durante tres meses lo eliminan, y después, con el mercado, hacen lo que quieren. Eso constituye precios predatorios, una conducta anticompetitiva.
\end{example}

\section{Dumping y derecho antidumping}

Relacionado con los precios predatorios, el docente plantea qué ocurre cuando ingresa un producto de otro país (generalmente de China, país muy grande que produce cantidades siderales, o de otros países) a un precio muy bajo, incluso más bajo que el que tiene en su país de origen, para tratar de conquistar el mercado. Los productores nacionales cuentan con una defensa: la ley o el \textbf{derecho antidumping}. El docente aclara que no quiere confundirlo con lo anterior, pero que conviene relacionarlos.

\begin{definition}{Dumping y derecho antidumping}
\textbf{Dumping:} cuando un producto entra al país a un precio más bajo del que se vende en el país de origen. Allí, según el docente, claramente hay algún tipo de subsidio estatal o de interés en romper el mercado.

\textbf{Derecho antidumping:} el país fija un derecho que compensa ese diferencial de precio para que el productor nacional no se vea desplazado. Si las diferencias son del 20\,\%, se fija un derecho del 20\,\% para que el productor nacional compita en igualdad de condiciones.

\textbf{Prueba:} no basta con que el producto sea muy barato. Hay que demostrar que el precio al cual ingresa el bien en la Argentina es más bajo que el precio al que lo vende en su país de origen el productor: allí está la inconducta comercial.
\end{definition}

\begin{table}[H]
\centering
\renewcommand{\arraystretch}{1.25}
\begin{tabularx}{\textwidth}{
    >{\raggedright\arraybackslash}p{0.2\textwidth}
    >{\raggedright\arraybackslash}X
    >{\raggedright\arraybackslash}X
}
\toprule
 & \textbf{Precios predatorios} & \textbf{Dumping} \\
\midrule
\textbf{Qué es} & Vender por debajo de los propios costos, o mantener un precio artificialmente bajo, para excluir competidores & Ingresar un producto a un precio inferior al de su país de origen \\
\textbf{Marco} & Defensa de la competencia: es una inconducta, una violación de la ley & Norma de comercio exterior \\
\textbf{Respuesta} & Procedimiento administrativo sancionatorio y eventual sanción & Derecho antidumping que compensa la diferencia de precio \\
\textbf{Qué hay que demostrar} & Precio bajo los costos, sostenido en el tiempo, que impide competir & Que el precio de ingreso es menor al del país de origen \\
\bottomrule
\end{tabularx}
\end{table}

El dumping es una norma de comercio exterior; la inconducta de competencia es mantener por mucho tiempo un precio predatorio, por debajo de los propios costos.

\section{Las dos herramientas de la Autoridad y factores de análisis de la posición dominante}

\subsection{La fase sancionatoria y la raíz preventiva}

La \textbf{fase sancionatoria} existe cuando se demuestra, a lo largo de un procedimiento administrativo sancionatorio, si hay o no una inconducta y una violación de la ley. Si al demandado o acusado se le demuestra la realización de una conducta anticompetitiva, se lo puede sancionar. Es la otra gran herramienta de la Autoridad de Defensa de la Competencia.

La primera herramienta es de raíz preventiva, a fin de no generar posiciones dominantes en el mercado: ya sea una posición dominante dentro del mismo mercado, o una posición dominante porque la empresa está en un mercado y además maneja el insumo que es clave de ese mercado, como en el caso de Coca-Cola que compra la embotelladora. Otro ejemplo, que el docente reconoce como absurdo: una empresa de fideos que compra la mayor cantidad de harina del país para manejar el principal insumo de su producto.

\subsection{Factores a tener en cuenta}

\begin{important}
\textbf{Cuatro factores para analizar la competencia y la posición dominante}
\begin{enumerate}
    \item ¿En qué mercado se encuentra la empresa?
    \item ¿Hay barreras de ingreso o egreso de ese mercado? ¿Cuán regulado está ese mercado?
    \item ¿Cuántos competidores hay y qué porcentaje tiene cada uno?
    \item La \textbf{sustitución}: puede haber un mercado relevante pero con muchos reemplazos. Si el mercado es relevante pero hay muchos productos de otros mercados que pueden reemplazarlo, «la sustitución debilita también esta situación»: puede estar muy concentrado y aun así haber muchos productos que sustituyan al que se está concentrando.
\end{enumerate}
\end{important}

\section{El procedimiento sancionatorio y la jerarquía constitucional}

\subsection{Un juzgado penal de instrucción}

La fase sancionatoria actúa como una suerte de \textbf{juzgado penal de instrucción}. El tribunal puede recibir una denuncia o hacer una investigación y, a raíz de ella, inicia el procedimiento para determinar si corresponde sancionar la violación a la ley: debe realizarse un procedimiento para establecer si hay o no una conducta anticompetitiva.

En el caso de las cementeras, las empresas coordinaban los precios, lo que se veía en los mercados, y además se demostró que se reunían en la cámara empresarial; se demostró así la idea de una connivencia.

% TODO: contenido incompleto o ambiguo en las notas originales
% (El apunte aclara que la expresión previa a «porque coordinaban los precios» resulta poco clara en el audio y se transcribió como «con las manos en la masa», coherente con el sentido de la frase.)

\subsection{Rango constitucional y necesidad de la política de competencia}

La ley antitrust, la ley de defensa de la competencia, tiene también \textbf{jerarquía constitucional}, y es una necesidad que todas las políticas estén orientadas a evitar concentraciones desmedidas y a fomentar una verdadera competencia. El mercado no es perfecto: puede haber quienes se aprovechan de una situación que quizás, al inicio, fue virtuosa, pero que, luego de mucho tiempo, cuando ya son dueños del mercado y no están expuestos a una real competencia, pueden generar distorsiones y abusar de esa posición.

Es clave que exista un tribunal que, si bien está inserto en la administración, sea un tribunal independiente: conforme a la ley, tiene inamovilidad en los cargos para que la ley no se use de forma arbitraria ni parcial, ni para castigar a enemigos y bendecir amigos.

\section{Cuadro Cornell del capítulo}

\begin{longtable}{@{}>{\raggedright\arraybackslash}p{\cornq}!{\vrule width 0.4pt}>{\raggedright\arraybackslash}p{\corna}@{}}
\toprule
\textbf{Preguntas / palabras clave} & \textbf{Notas / respuestas} \\
\midrule
\endfirsthead
\toprule
\textbf{Preguntas / palabras clave} & \textbf{Notas / respuestas} \\
\midrule
\endhead
\midrule
\multicolumn{2}{r}{\small\itshape continúa en la página siguiente} \\
\endfoot
\bottomrule
\endlastfoot

\textbf{¿Qué conductas son anticompetitivas?} &
\textbf{Venta atada}, \textbf{cartelización} (cementeras, prepagas), \textbf{reparto de zonas}, restricción de cantidades, bloqueo de competidores y \textbf{precios predatorios} (por debajo del costo para fundir al rival o sostenidos en el tiempo; debe probarse). \\
\addlinespace[0.6em]

\textbf{¿En qué se diferencia el dumping de los precios predatorios?} &
El \textbf{dumping} es ingresar un producto a un precio menor que en su país de origen; se responde con un \textbf{derecho antidumping} que compensa la diferencia (norma de comercio exterior). El precio predatorio es una inconducta de competencia que se sanciona en un procedimiento administrativo. \\
\addlinespace[0.6em]

\textbf{¿Cómo actúa la vertiente sancionatoria y qué jerarquía tiene la ley?} &
Actúa como una suerte de \textbf{juzgado penal de instrucción}: recibe una denuncia o investiga, e inicia un procedimiento para determinar si hubo conducta anticompetitiva. La ley tiene \textbf{jerarquía constitucional}, y el tribunal es independiente, con inamovilidad en los cargos, para que la ley no se use arbitrariamente. \\

\midrule
\multicolumn{2}{@{}p{\dimexpr\cornq+\corna+2\tabcolsep\relax}@{}}{%
\textbf{Resumen:} La vertiente sancionatoria castiga el abuso de posición dominante y las conductas anticompetitivas (cartelización, venta atada, reparto de zonas, restricción de cantidades, precios predatorios), previa demostración en un procedimiento administrativo. El dumping, en cambio, se rige por una norma de comercio exterior y se responde con un derecho antidumping.} \\
\end{longtable}

% ============================================================
\chapter[Garantías del control administrativo]{Garantías del control administrativo y cierre}
% ============================================================

\section{Requisitos de la Corte Suprema para las sanciones administrativas}

La jurisprudencia de la Corte Suprema se refiere, sobre todo, a los entes administrativos que aplican sanciones y resuelven situaciones entre particulares como si fueran, de algún modo, jueces (no lo son). Todos los organismos administrativos que imponen sanciones, o que tienen poder de decisión sobre conflictos entre particulares, deben cumplir una serie de requisitos para poder emitir o fijar sanciones válidamente.

\begin{important}
\textbf{Requisitos para que un organismo administrativo pueda sancionar válidamente}
\begin{enumerate}
    \item \textbf{Potestad sancionatoria fijada por ley.} No puede ser un decreto ni una norma reglamentaria: tiene que haber una ley formal que habilite la competencia sancionatoria.
    \item \textbf{Idoneidad técnica justificada y ámbito específico.} Tiene que ser razonable que el tribunal tenga esa facultad sancionatoria porque cuenta con capacitación y hay un ámbito específico. No puede haber grandes ámbitos generales de sanciones; tiene que haber ámbitos concretos fijados por la norma, porque si no se estarían creando tribunales de facto.
    \item \textbf{Independencia.} Condiciones de inamovilidad en los cargos: su conducta no puede estar a merced de la decisión del gobierno de turno.
    \item \textbf{Imparcialidad.} Deben ser independientes y asegurar idoneidad técnica, pero también imparcialidad.
    \item \textbf{Control judicial suficiente.} Cualquier decisión, sobre todo cualquier sanción, tiene que ser apelable.
\end{enumerate}
\end{important}

\begin{note}
El apunte indica que el docente no nombró los precedentes, y que esta doctrina se desarrolla, entre otros, en «Fernández Arias c/ Poggio» (1960) y «Ángel Estrada y Cía.» (2005), ambos de la Corte Suprema.
\end{note}

\subsection{El control judicial suficiente}

Está presente en todas estas normas: cualquier decisión que se tome, sobre todo cualquier sanción, tiene que ser apelable, es decir, debe existir un control judicial suficiente. Tiene que haber una vía judicial para impugnarla. No puede quedar con la última palabra un órgano administrativo; de lo contrario, no habría un Estado de derecho.

La decisión administrativa, que puede implicar, por ejemplo, rechazar una concentración económica o imponer una multa, tiene que ser revisable por la justicia. La justicia, si bien entiende que el organismo idóneo y especializado es el administrativo, va a revisar que se cumpla con el debido proceso, que no haya una decisión irrazonable o infundada y que se hayan respetado todas las garantías en el procedimiento administrativo sancionatorio. Hay una revisión judicial suficiente: puede abarcar tanto el derecho como los hechos. De este modo no queda ninguna restricción al control judicial y se asegura que las arbitrariedades puedan ser anuladas en la etapa ante el Poder Judicial.

\section{Cierre de la clase}

El docente cierra con un mensaje para retener los conceptos: analizar bien, de cada una de las dos leyes, quién es el organismo de aplicación y cuáles son los conceptos generales.

\begin{note}
\textbf{Tarea indicada para la próxima clase.} Para cada una de las dos leyes (consumidor y competencia), identificar quién es el organismo de aplicación y cuáles son los conceptos generales. El docente enviará material complementario.
\end{note}

\section{Cuadro Cornell del capítulo}

\begin{longtable}{@{}>{\raggedright\arraybackslash}p{\cornq}!{\vrule width 0.4pt}>{\raggedright\arraybackslash}p{\corna}@{}}
\toprule
\textbf{Preguntas / palabras clave} & \textbf{Notas / respuestas} \\
\midrule
\endfirsthead
\toprule
\textbf{Preguntas / palabras clave} & \textbf{Notas / respuestas} \\
\midrule
\endhead
\midrule
\multicolumn{2}{r}{\small\itshape continúa en la página siguiente} \\
\endfoot
\bottomrule
\endlastfoot

\textbf{¿Qué requisitos exige la Corte para que un organismo administrativo sancione?} &
(1) \textbf{Ley formal} que otorgue la potestad; (2) \textbf{idoneidad técnica} y ámbito específico; (3) \textbf{independencia} e inamovilidad; (4) \textbf{imparcialidad}; (5) \textbf{control judicial suficiente}, apelable y revisable en hechos y derecho, para que no quede la última palabra en la administración. \\
\addlinespace[0.6em]

\textbf{¿Por qué es necesario el control judicial suficiente?} &
Porque cualquier decisión, sobre todo cualquier sanción, debe ser apelable: si un órgano administrativo tuviera la última palabra no habría Estado de derecho. La justicia revisa el debido proceso, la razonabilidad y la fundamentación, tanto en el derecho como en los hechos, de modo que las arbitrariedades puedan ser anuladas. \\

\midrule
\multicolumn{2}{@{}p{\dimexpr\cornq+\corna+2\tabcolsep\relax}@{}}{%
\textbf{Resumen:} Todo organismo administrativo que sanciona o resuelve conflictos entre particulares debe contar con ley formal, idoneidad técnica en un ámbito específico, independencia, imparcialidad y control judicial suficiente. Las decisiones administrativas, como multas o rechazos de concentraciones, siempre deben poder ser revisadas por un juez.} \\
\end{longtable}

% ============================================================
\chapter{Síntesis final}
% ============================================================

La clase mostró cómo el Estado regula la economía desde dos ángulos complementarios y con jerarquía constitucional. La defensa del consumidor protege a la parte débil de cada relación de consumo mediante una ley de orden público (nulidad de cláusulas abusivas, información veraz, trato digno, solidaridad en la cadena, daño punitivo y acciones de clase promovidas por asociaciones registradas), mientras que la defensa de la competencia cuida el funcionamiento mismo del mercado.

Lo hace con una vertiente preventiva, que controla las concentraciones económicas y puede aprobarlas, condicionarlas o rechazarlas según el mercado relevante, las barreras de entrada, la cantidad de competidores y la sustitución, y con una vertiente sancionatoria, que castiga el abuso de posición dominante y las conductas anticompetitivas como la cartelización, la venta atada y los precios predatorios.

La autoridad de aplicación debe ser independiente (cinco miembros con acuerdo del Senado e inamovilidad) y, como todo organismo administrativo que sanciona, cumplir los requisitos de la Corte: ley formal, idoneidad técnica, independencia, imparcialidad y control judicial suficiente.

La idea que une todo es que ni el mercado ni la administración pueden quedar sin controles: el consumidor debe tener a quién reclamar, la competencia debe ser real y el poder de sancionar debe poder ser revisado por un juez.

\end{document}
